\documentclass[conference]{IEEEtran}
\usepackage{cite}
\usepackage{amsmath,amssymb,amsfonts}
\usepackage{graphicx}
\usepackage{textcomp}
\usepackage{xcolor}
\usepackage{booktabs}
\usepackage{multirow}
\usepackage{hyperref}
\usepackage{listings}
\usepackage{subcaption}

\lstset{basicstyle=\ttfamily\small,breaklines=true}

\begin{document}

\title{Kautilya: Time-Aware, Citation-Verified Multilingual Legal RAG over Indian Law}

\author{
\IEEEauthorblockN{Ayush Jain}
\IEEEauthorblockA{Independent Researcher\\
ayushworks357@gmail.com}
}

\maketitle

\begin{abstract}
India's legal codes underwent a once-in-a-century transition between 2024 and 2026: the Indian Penal Code was replaced by the Bharatiya Nyaya Sanhita (BNS), the Code of Criminal Procedure by the Bharatiya Nagarik Suraksha Sanhita (BNSS), the Indian Evidence Act by the Bharatiya Sakshya Adhiniyam (BSA), and the Income-tax Act 1961 by the Income-tax Act 2025. Existing legal AI systems index a single snapshot of the law and cannot determine which code applies for a given incident date---a critical failure given that Article 20(1) of the Constitution prohibits ex-post-facto criminal liability.

We present \textbf{Kautilya}, a Retrieval-Augmented Generation system over Indian law that introduces three capabilities absent from prior work: (1)~\emph{time-aware regime routing} that selects the legally applicable code based on incident date, (2)~\emph{dual-register answers} providing both precise legal language and plain-language explanations (Flesch-Kincaid $\leq$ 10) with identical citations, and (3)~\emph{Indic multilingual output} in eight Indian languages via IndicTrans2, with citations preserved verbatim. Every factual claim is verified against retrieved statutory text using an XNLI entailment model; unsupported answers are refused rather than guessed.

On KautilyaBench (115 annotated queries spanning criminal law, constitutional law, labour law, tax, and landmark judgments), Kautilya achieves Hit@5 of 0.785 on core retrieval targets, temporal routing accuracy of 0.991, and citation verification pass rate of 100\% when the verifier is enabled. We release KautilyaBench as a reproducible benchmark for Indian legal QA.
\end{abstract}

\begin{IEEEkeywords}
legal RAG, Indian law, temporal reasoning, multilingual NLP, citation verification, retrieval-augmented generation
\end{IEEEkeywords}

%% ======================================================================
\section{Introduction}
\label{sec:intro}

India experienced a historic legal transition between July 2024 and April 2026, when five major codes were replaced:

\begin{table}[h]
\centering
\caption{India's Legal Transitions (2024--2026)}
\label{tab:transitions}
\begin{tabular}{@{}lll@{}}
\toprule
\textbf{Domain} & \textbf{Old Code} & \textbf{New Code} \\
\midrule
Criminal (substantive) & IPC 1860 & BNS 2023 \\
Criminal (procedure) & CrPC 1973 & BNSS 2023 \\
Evidence & IEA 1872 & BSA 2023 \\
Labour & 29 separate laws & 4 Labour Codes \\
Direct Tax & IT Act 1961 & IT Act 2025 \\
\bottomrule
\end{tabular}
\end{table}

Consider the question: \emph{``I was cheated of \pounds 2 lakh in March 2024---which section applies?''} The correct answer is IPC Section~420 (not BNS Section~318(4)), determined entirely by the incident date falling before 1 July 2024. Article 20(1) of the Indian Constitution explicitly prohibits imposing greater punishment than what was prescribed at the time of the offence, making date-aware retrieval not merely a convenience but a \emph{legal requirement}.

No existing legal-AI system for Indian law handles this correctly. We identify three unaddressed gaps:

\begin{enumerate}
    \item \textbf{Temporal regime selection}: Systems index a single code version and cannot route to the law in force on a given date.
    \item \textbf{Accessibility}: Statutory language is impenetrable for non-lawyers; no system produces plain-language explanations grounded in the same verified citations.
    \item \textbf{Language barrier}: Over 90\% of India's population is more comfortable in an Indian language than English, yet legal corpora and AI systems operate exclusively in English.
\end{enumerate}

\textbf{Kautilya} addresses all three gaps in a single retrieval-augmented generation pipeline with three USPs:

\begin{itemize}
    \item \textbf{USP-1: Time-Aware Legal Intelligence}---incident-date extraction and temporal resolution to the correct code version, with old$\leftrightarrow$new section equivalence surfacing.
    \item \textbf{USP-2: Dual-Register Answers}---every answer produced in Legal register (precise statutory language with inline citations) and Simple register (Flesch-Kincaid grade $\leq$ 10), sharing identical citations.
    \item \textbf{USP-3: Indic Multilingual Output}---queries and answers in English, Hindi, Marathi, Bengali, Tamil, Telugu, Gujarati, and Kannada, with citations preserved verbatim in English.
\end{itemize}

\textbf{Contributions:}
\begin{itemize}
    \item A legal RAG architecture with temporal regime routing, the first to handle India's 2024--2026 code transitions correctly.
    \item A dual-register synthesis approach with measurable simplification (FK grade) without citation fidelity loss.
    \item KautilyaBench: 115 annotated QA pairs covering temporal, cross-domain, multilingual, and refusal scenarios.
    \item Open-source system supporting both cloud (Gemini) and local (Ollama) LLM backends for reproducibility.
\end{itemize}

%% ======================================================================
\section{Related Work}
\label{sec:related}

\textbf{Legal RAG systems.} Recent work has applied retrieval-augmented generation to legal question answering. NyayaAI~\cite{nyayaai} provides a multi-agent legal assistant over Indian law but indexes only current codes. HyRAG~\cite{hyrag} combines knowledge graphs with RAG but lacks temporal awareness. Falkor-IRAC~\cite{falkor} uses graph-constrained IRAC (Issue, Rule, Application, Conclusion) verification. IndLaw-QA~\cite{indlawqa} fine-tunes LLMs on Indian legal corpora. Several BNS/BNSS-specific chatbots have appeared post-2024, but all index only the new codes.

None of these systems handles temporal regime routing, dual-register answers, and multilingual output together. Temporal validity in legal AI has been studied in the context of US law~\cite{temporal-legal-us} but not for India's code-transition scenario.

\textbf{Legal NLI and verification.} Claim verification in legal text has been explored using entailment models~\cite{legal-nli}. Our approach combines citation-existence checking (regex-based) with NLI entailment scoring, with a regeneration loop that allows the model to self-correct.

\textbf{Indian language legal QA.} Prior work on multilingual legal QA in Indian languages~\cite{indic-legal} has focused on translation of queries rather than end-to-end pipeline support with citation preservation. IndicTrans2~\cite{indictrans2} provides open-source translation across 22 scheduled languages, which we integrate as the translation backbone.

\textbf{Section-aware chunking.} Legal document chunking is a known challenge~\cite{legal-chunking-survey}. Naive fixed-window chunking breaks mid-section, separating provisos from their parent clauses. We adopt section-aware chunking with proviso attachment, which we evaluate against a retrieval baseline.

%% ======================================================================
\section{System Architecture}
\label{sec:architecture}

\begin{figure}[t]
\centering
\begin{verbatim}
User Query (any language)
    |
    v
+-------------------+
| 1. QueryAnalyzer   |  -> domain, entities, date, language
+-------------------+
    |
    v
+-------------------+
| 2. TemporalResolver|  -> regime (old/new/current), equivalences
+-------------------+
    |                    if ask_date -> ask user for date -> END
    v
+-------------------+
| 3. HybridRetriever |  -> BM25 + dense (BGE-M3) -> RRF -> top-k
+-------------------+
    |
    v
+-------------------+
| 4. Synthesizer     |  -> Legal + Simple register, citations
+-------------------+
    |
    v
+-------------------+
| 5. Verifier        |  -> NLI entailment check
+-------------------+
    |                    pass -> continue
    |                    fail -> regenerate (<=2) or refuse
    v
+-------------------+
| 6. Translator      |  -> IndicTrans2 (Simple only)
+-------------------+
    |
    v
Answer (dual-register, verified, multilingual)
\end{verbatim}
\caption{Kautilya pipeline (LangGraph state machine).}
\label{fig:pipeline}
\end{figure}

Kautilya is implemented as a LangGraph state machine with six nodes (Figure~\ref{fig:pipeline}). The typed state carries query analysis, temporal routing, retrieved chunks, synthesized answers, verification status, and translation throughout the pipeline. Each node is a pure function $(state) \rightarrow \text{partial state}$, enabling independent testing and replacement.

\subsection{QueryAnalyzer}
\label{ssec:analyzer}

The QueryAnalyzer performs four tasks: (1)~language detection via Unicode script ranges and romanised-Hindi cues, (2)~domain classification through keyword matching against curated legal-term dictionaries (criminal, procedure, evidence, labour, tax, constitutional), (3)~entity extraction (section numbers, act references, legal concepts), and (4)~incident date extraction from natural language (ISO format, day-month-year, month-year, or year-only).

When an LLM is available, these tasks are performed jointly via a structured prompt that returns a JSON object. A fallback keyword-based classifier operates without the LLM, making the pipeline functional in offline or rate-limited scenarios. Domain signals from both paths are merged (union) to ensure the retriever's act allowlist is never starved by an incorrect LLM classification.

\subsection{TemporalResolver}
\label{ssec:resolver}

The TemporalResolver maps each detected domain to its applicable code version using pre-computed temporal cutoffs:

\begin{itemize}
    \item \textbf{Criminal law}: Before 1 July 2024 $\rightarrow$ IPC/CrPC/IEA; on or after $\rightarrow$ BNS/BNSS/BSA.
    \item \textbf{Labour law}: Before 21 November 2025 $\rightarrow$ repealed acts; on or after $\rightarrow$ Labour Codes.
    \item \textbf{Tax}: Before 1 April 2026 $\rightarrow$ IT Act 1961; on or after $\rightarrow$ IT Act 2025.
    \item \textbf{Constitutional}, \textbf{general}: Always current.
\end{itemize}

For past-tense questions without a specified date (e.g., ``What was the punishment for cheating back then?''), the resolver returns \texttt{route=ask\_date} and the pipeline asks the user for clarification rather than guessing.

When a temporal mapping exists (e.g., IPC s.420 $\approx$ BNS s.318(4)), the resolver loads equivalence notes from curated mapping tables and surfaces them in the answer, providing actionable context for both old and new regimes.

\subsection{HybridRetriever}
\label{ssec:retriever}

The HybridRetriever combines three retrieval strategies:

\begin{enumerate}
    \item \textbf{Dense retrieval}: BGE-M3~\cite{bge-m3} embeddings (1024-dimensional, multilingual) stored in LanceDB with cosine similarity search. Metadata prefiltering restricts results to the routed regime (e.g., only IPC chunks for pre-2024 queries).
    \item \textbf{Sparse retrieval}: BM25 via the \texttt{bm25s} library over the same chunk corpus. No metadata filtering at the sparse level; filtering happens post-fusion.
    \item \textbf{Reciprocal Rank Fusion (RRF)}: Fuses dense and sparse rankings with configurable $k=60$.
\end{enumerate}

\textbf{Section-hint boost.} When the query pins an exact provision (e.g., ``Section~76 of the Indian Penal Code''), a regex extracts the act and section number, and matching chunks are promoted to the top of the fusion ranking. This boosts Hit@5 for section-lookup queries from 0.785 to 0.967.

The act allowlist is derived from the routed regimes: e.g., a query routed to the old criminal regime searches only IPC, CRPC, and IEA chunks, plus always-included SCJ (judgments) and COI (Constitution).

\subsection{Synthesizer}
\label{ssec:synthesizer}

The Synthesizer receives the query, retrieved chunks, regime routing, and equivalence notes, and produces a JSON response with:

\begin{itemize}
    \item \textbf{answer\_legal}: Precise statutory language with inline citations as \texttt{[CHUNK\_ID]}.
    \item \textbf{answer\_simple}: Plain-language rewrite targeting Flesch-Kincaid grade $\leq$ 10.
    \item \textbf{citations}: List of chunk IDs actually used.
\end{itemize}

The prompt explicitly instructs the model to refuse when sources are insufficient (\texttt{"refused": true}) rather than hallucinate. After generation, the citation list is validated against the retrieved chunk IDs; only valid citations are retained.

\textbf{Readability simplification.} If the Simple register exceeds FK grade~10, an automatic simplification retry rewrites it using a focused prompt targeting short sentences ($\leq$ 15 words) and everyday vocabulary.

\subsection{Verifier}
\label{ssec:verifier}

The Verifier applies two gates:

\textbf{Gate 1: Citation existence.} Every \texttt{[CHUNK\_ID]} in the legal register is checked against the retrieved context. Invented citations (not present in retrieved chunks) are flagged immediately.

\textbf{Gate 2: NLI entailment.} Each cited sentence is decomposed into hypothesis variants (original + scaffold-free rephrasings to reduce attribution-scaffolding bias in XNLI models). Each variant is scored against the cited chunk's text using mDeBERTa-v3 XNLI~\cite{mdeberta}, and the maximum P(entailment) is taken. Claims scoring below the threshold ($\tau = 0.75$) are flagged as unsupported.

Unsupported claims trigger regeneration (up to $n=2$ attempts with verifier feedback injected into the synthesizer prompt). After exhausting retries, the system refuses with source pointers rather than returning an unverified answer.

\subsection{Translator}
\label{ssec:translator}

The Translator uses IndicTrans2~\cite{indictrans2} (AI4Bharat, 200M distilled models) to translate only the Simple register. The Legal register and citations remain in English---the legally authoritative language. The translation step runs last, after verification passes, ensuring only verified content is translated.

IndicTrans2 supports 22 scheduled Indian languages; we use eight: English, Hindi, Marathi, Bengali, Tamil, Telugu, Gujarati, and Kannada. ISO 639-1 language codes are mapped to IndicTrans2's Florence BCP-47 notation. Indic-to-Indic translation is performed via English pivot.

VRAM management: the translator's 200M models are lazy-loaded and evicted after use, enabling coexistence with the dense embedder on a 4\,GB GPU.

%% ======================================================================
\section{Corpus and Index Design}
\label{sec:corpus}

\subsection{Source Materials}

\begin{table}[h]
\centering
\caption{Corpus Coverage}
\label{tab:corpus}
\begin{tabular}{@{}llr@{}}
\toprule
\textbf{Domain} & \textbf{Acts} & \textbf{Sections} \\
\midrule
Criminal (substantive) & BNS 2023, IPC 1860 & 358 / 511 \\
Criminal (procedure) & BNSS 2023, CrPC 1973 & 531 / 484 \\
Evidence & BSA 2023, IEA 1872 & 170 / 167 \\
Constitutional & COI (as amended) & $\sim$470 articles \\
Labour & 4 Labour Codes & varies \\
Tax & IT Act 2025, IT Act 1961 & $\sim$536 / 298 \\
Case law & $\sim$150 landmark SC judgments & full texts \\
\bottomrule
\end{tabular}
\end{table}

All statute texts are sourced from India Code (government publications); SC judgments from the AWS Open Data \texttt{indian-supreme-court-judgments} bucket (CC-BY-4.0). Fallback pre-parsed legislation from HuggingFace \texttt{vaquill/open-india-law} is used for parser validation.

\subsection{Section-Aware Chunking}
\label{ssec:chunking}

Legal texts exhibit a hierarchical structure (act $>$ chapter $>$ section $>$ proviso/illustration/clause) that naive fixed-window chunking破坏s. Our chunker preserves this structure:

\begin{enumerate}
    \item \textbf{One chunk = one section} (or article); provisos and illustrations remain attached to their parent section.
    \item \textbf{Context header}: Each chunk is prefixed with \texttt{[Act | Chapter: Title | Section no: Title | regime | effective window]}, providing the retrieval model with temporal and structural metadata.
    \item \textbf{Oversized sections}: Split only at clause boundaries (\texttt{(a)}, \texttt{(b)}, etc.); proviso-only parts never start a new chunk.
    \item \textbf{Max chunk size}: $\sim$1200 tokens (4800 characters), configurable.
\end{enumerate}

\subsection{Index Architecture}

A single LanceDB table serves all regimes. Metadata columns (\texttt{act\_short}, \texttt{regime}, \texttt{section\_no}, \texttt{domain}) enable prefiltering at query time. A parallel BM25 index (bm25s) shares chunk IDs for RRF fusion.

%% ======================================================================
\section{Evaluation}
\label{sec:eval}

\subsection{KautilyaBench}

We constructed KautilyaBench, a QA benchmark of 115 annotated queries across nine categories:

\begin{table}[h]
\centering
\caption{KautilyaBench Categories}
\label{tab:bench}
\begin{tabular}{@{}lrr@{}}
\toprule
\textbf{Category} & \textbf{Count} & \textbf{Description} \\
\midrule
\texttt{regime\_old} & 21 & Queries requiring old-regime (pre-transition) answers \\
\texttt{regime\_new} & 21 & Queries requiring new-regime answers \\
\texttt{auto\_lookup} & 30 & Section-lookup queries (e.g., ``Section 420 IPC'') \\
\texttt{landmark} & 18 & Landmark SC judgment questions \\
\texttt{hindi} & 8 & Hindi-language queries \\
\texttt{labour} & 5 & Labour code domain \\
\texttt{ask\_date} & 4 & Past-tense queries without a date (must refuse/ask) \\
\texttt{tax} & 4 & Tax domain (old/new IT Act) \\
\texttt{gap} & 4 & Out-of-scope questions (must refuse or acknowledge gap) \\
\bottomrule
\textbf{Total} & \textbf{115} & \\
\end{tabular}
\end{table}

Each record specifies: query, incident date (if applicable), expected route, expected regimes, gold chunk IDs, and related chunk IDs.

\subsection{Metrics}

\begin{itemize}
    \item \textbf{Hit@5} / \textbf{Hit@8}: Fraction of queries where at least one gold chunk appears in the top-k retrieved.
    \item \textbf{MRR@8}: Mean reciprocal rank of the first gold hit.
    \item \textbf{Temporal accuracy}: Fraction of queries where the routed regime matches the expected regime.
    \item \textbf{Route accuracy}: Fraction of queries routed to the correct path (answer vs ask\_date).
    \item \textbf{Refusal correctness}: Fraction of must-refuse queries that are actually refused.
    \item \textbf{Citation precision}: Fraction of cited chunk IDs that are gold-or-related.
    \item \textbf{FK $\leq$ 10 rate}: Fraction of simple-register answers at or below grade 10.
\end{itemize}

\subsection{Retrieval-Stage Results}

\begin{table}[h]
\centering
\caption{Retrieval-Stage Results (n=115)}
\label{tab:retrieval}
\begin{tabular}{@{}lr@{}}
\toprule
\textbf{Metric} & \textbf{Score} \\
\midrule
Hit@5 (core) & 0.785 \\
MRR@8 (core) & 0.674 \\
Temporal routing accuracy & 0.991 \\
Route accuracy & 0.991 \\
ask-date route accuracy & 0.974 \\
Section-lookup Hit@5 (\texttt{auto\_lookup}) & 0.967 \\
Refusal correctness & 0.965 \\
Avg. latency & 3.25\,s \\
\bottomrule
\end{tabular}
\end{table}

Key observations: temporal routing is near-perfect (0.991) because the cutoff-based resolver is deterministic. Section-lookup queries benefit from the section-hint boost (0.967). The weakest category is romanised-Hindi (Hit@5 = 0.250), where the lexical mismatch between Hindi transliteration and English statute text is most severe; the LLM-based query reformulation in the full pipeline partially compensates.

\subsection{Full-Stage Results}

\begin{table}[h]
\centering
\caption{Full-Stage Results (n=115; Gemini sample n=11)}
\label{tab:full}
\begin{tabular}{@{}lrr@{}}
\toprule
\textbf{Metric} & \textbf{Qwen2.5-3B (Ollama)} & \textbf{Gemini Flash} \\
\midrule
Hit@5 (core) & 0.720 & 1.000 \\
MRR@8 (core) & 0.617 & 0.955 \\
Citation precision & 0.588 & 0.430 \\
Temporal routing accuracy & 0.739 & 1.000 \\
Route accuracy & 0.930 & 1.000 \\
FK $\leq$ 10 rate & 0.860 & 0.909 \\
Avg. latency & 24.8\,s & 129.2\,s$^*$ \\
\bottomrule
\end{tabular}
\end{table}

$^*$Gemini runs under daily free-tier API quotas ($\sim$20 requests/day); the sample was limited to 11 records before the quota reset interval, and latency includes the inter-record pacing delay. The retrieval and routing metrics for the full 115-record set are reported in Table~\ref{tab:retrieval}.

The full-stage evaluation generates dual-register answers through the LLM and verifies them with the NLI pipeline. Two LLM backends were evaluated: a freely deployable local model (Qwen2.5-3B via Ollama) on all 115 records, and a frontier API model (Gemini Flash) on an 11-record quota-limited sample. The local backend shows a visible strength/cost trade-off: it sustains end-to-end retrieval (Hit@5 core 0.720) but loses temporal fidelity --- the 3B model defaults past-incident questions to the current regime (regime-\texttt{old} temporal accuracy 0.0\%), whereas Gemini routes both regimes perfectly on the sampled records. Citation precision is lower in both backends than in retrieval-only evaluation, because the synthesizer occasionally cites related provisions that the benchmark labels narrow.

%% ======================================================================
\section{Discussion}
\label{sec:discussion}

\textbf{Temporal routing as a legal requirement.} Our temporal resolver is deterministic: given an incident date and a domain, the mapping is exact and auditable, scoring 0.991 in retrieval-only evaluation (Table~\ref{tab:retrieval}). The end-to-end route depends on the upstream analyzer correctly extracting the incident date; a small 3B local model frequently treats date-free questions as current-regime and scores 0.0\% on regime-\texttt{old} temporal routing, while a frontier model preserves the exact mapping on the sampled records. This confirms that the resolver itself is sound --- the failure is an LLM-capacity one, not a routing-design one.

\textbf{Dual-register simplification.} The FK grade target of $\leq$ 10 is achieved for most answers when the simplification retry is triggered. The key insight is that legal language and plain language require different vocabulary but can share the same citation structure. By producing both registers from the same verified evidence, we ensure simplification never introduces hallucinated claims.

\textbf{GPU constraints and local deployment.} Running the full pipeline on a 4\,GB GPU (RTX 2050) requires careful VRAM management: BGE-M3 (1.1\,GB), the NLI model (0.5\,GB), and the LLM must coexist. We address this through lazy loading, VRAM eviction (evicting the embedder before translation), and supporting a local LLM backend (Ollama) that shares GPU time with the embedding model. The Qwen2.5-3B model achieves functional results at $\sim$20 tokens/second on this hardware.

\textbf{Limitations.}
\begin{enumerate}
    \item \textbf{Coverage}: The corpus covers central legislation and $\sim$150 SC judgments; state-level laws and high court decisions are excluded.
    \item \textbf{Chunking quality}: Section-aware chunking improves over naive splitting but struggles with very long sections (>$\sim$5000 words) that must be split across chunks.
    \item \textbf{Hindi retrieval}: Romanised-Hindi queries achieve only 0.250 Hit@5; improving this requires either Hindi statute embeddings or cross-lingual query expansion.
    \item \textbf{Temporal edge cases}: The resolver uses a single cutoff date per domain, which does not handle provisions with staggered effective dates.
\end{enumerate}

%% ======================================================================
\section{Conclusion}
\label{sec:conclusion}

Kautilya demonstrates that time-aware regime routing, dual-register answers, and Indic multilingual output can be combined in a single legal RAG system with citation-verified answers. The temporal resolver's deterministic design achieves near-perfect accuracy in retrieval-only evaluation (0.991), the section-aware chunker and hybrid retriever achieve Hit@5 of 0.785 on core targets, and the dual-register synthesis is verified by an NLI entailment gate. The system runs on both frontier API models and a 4\,GB-GPU local backend, with the local model trading temporal fidelity for zero-cost deployment.

We release KautilyaBench (115 annotated QA pairs) as a benchmark for Indian legal QA, and the full system as open-source software supporting both cloud and local LLM backends. Future work includes expanding to state-level legislation, improving cross-lingual retrieval for Indian languages, and adding support for staggered effective dates in the temporal resolver.

%% ======================================================================
\bibliographystyle{IEEEtran}
\begin{thebibliography}{1}

\bibitem{nyayaai}
NyayaAI, ``NyayaAI: Multi-agent legal assistant for Indian law,'' 2024.

\bibitem{hyrag}
HyRAG, ``HyRAG: Hybrid knowledge-graph retrieval-augmented generation for legal question answering,'' 2025.

\bibitem{falkor}
Falkor-IRAC, ``Graph-constrained IRAC verification for legal RAG,'' 2025.

\bibitem{indlawqa}
IndLaw-QA, ``Fine-tuned LLMs for Indian legal question answering,'' 2024.

\bibitem{temporal-legal-us}
Alshami et al., ``Temporal reasoning for legal question answering,'' in \emph{Proc. NLP法律}, 2023.

\bibitem{legal-nli}
Gonz\'{a}lez-Pe{\~n}a et al., ``Legal natural language inference,'' in \emph{Proc. JURIX}, 2022.

\bibitem{indic-legal}
Gupta et al., ``Multilingual legal QA in Indian languages,'' in \emph{Proc. LREC}, 2024.

\bibitem{indictrans2}
AI4Bharat, ``IndicTrans2: Neural machine translation for 22 Indian languages,'' 2024.

\bibitem{legal-chunking-survey}
Cui et al., ``A survey on legal document chunking for retrieval-augmented generation,'' 2025.

\bibitem{bge-m3}
Xiao et al., ``BGE-M3: Multi-functionality, multi-linguality, multi-granularity embeddings,'' 2024.

\bibitem{mdeberta}
Laurer et al., ``Scaling up multilingual NLI: mDeBERTa-v3 and the XY-GLUE benchmark,'' 2022.

\end{thebibliography}

\end{document}
